% Flowchart Template — ISO 5807 standard
% Copy this file as a starting point for every flowchart figure.
% Do NOT change styles or preamble. Layout constants may be adjusted
% in the CONTENT section if the default spacing is too tight or loose.
%
% Usage:
%   1. Copy to figures/fig_<name>/fig_<name>.tex
%   2. Fill in nodes (Column C, L1, R1) and arrows
%   3. Adjust layout constants if needed (move \def lines into CONTENT)
%   4. Compile: pdflatex fig_<name>.tex
%
\documentclass[border=8pt]{standalone}
\usepackage{tikz}
\usetikzlibrary{shapes.geometric, arrows.meta, positioning, calc}

\begin{document}

% =================================================================
% STYLES — do not modify
% =================================================================
\tikzset{
  % --- Node styles (ISO 5807) ---
  startstop/.style={
    stadium, minimum width=3.2cm, minimum height=1cm,
    text centered, draw=black, fill=blue!15,
    font=\sffamily
  },
  process/.style={
    rectangle, minimum width=3.6cm, minimum height=1cm,
    text centered, draw=black, fill=orange!12,
    font=\sffamily
  },
  decision/.style={
    diamond, aspect=2.2,
    minimum width=3.2cm, minimum height=1cm,
    text centered, draw=black, fill=green!12,
    font=\sffamily, inner sep=1pt
  },
  io/.style={
    trapezium, trapezium left angle=70, trapezium right angle=110,
    minimum width=3.2cm, minimum height=1cm,
    text centered, draw=black, fill=yellow!15,
    font=\sffamily
  },
  predef/.style={
    rectangle, double, double distance=2pt,
    minimum width=3.6cm, minimum height=1cm,
    text centered, draw=black, fill=violet!10,
    font=\sffamily
  },
  % --- Arrow style ---
  arrow/.style={
    thick, -{Stealth[length=6pt]}
  },
  % --- Edge label style ---
  lbl/.style={
    font=\sffamily\small, fill=white, inner sep=2pt
  },
}

% =================================================================
% LAYOUT CONSTANTS
% =================================================================
% Vertical spacing between center-column nodes
\def\vgapSmall{1.5cm}   % process → process, process → decision
\def\vgapLarge{2.2cm}    % decision → next node (diamond needs more room)
% Horizontal distance from center to side nodes (L1/R1)
\def\sideGap{3.5cm}
% Feedback loop clearance beyond outermost side node edge
\def\loopClearance{2.0cm}

% =================================================================
% CONTENT — replace the example below with your diagram
% =================================================================
\begin{tikzpicture}[node distance=1.8cm]

  % --- Column C (center spine) ---
  \node (start)  [startstop]                          {Start};
  \node (proc1)  [process, below=\vgapSmall of start] {Process 1};
  \node (dec1)   [decision, below=\vgapLarge of proc1]{Condition?};
  \node (proc2)  [process, below=\vgapLarge of dec1]  {Process 2};
  \node (stop)   [startstop, below=\vgapSmall of proc2]{End};

  % --- Column L1 (left side exits) ---
  % \node (retA) [startstop, left=\sideGap of dec1] {Return A};

  % --- Column R1 (right side exits) ---
  % \node (retB) [startstop, right=\sideGap of dec1] {Return B};

  % --- Main spine arrows ---
  \draw [arrow] (start) -- (proc1);
  \draw [arrow] (proc1) -- (dec1);
  \draw [arrow] (dec1)  -- node[lbl, right] {Yes} (proc2);
  \draw [arrow] (proc2) -- (stop);

  % --- Side exit arrows ---
  % \draw [arrow] (dec1) -- node[lbl, above] {No} (retA);

  % --- Feedback loops ---
  % Two-part algorithm (see flowchart_standard.md for full explanation):
  %
  % Part 1: Compute safe L2/R2 x from ALL side nodes on each side.
  %
  % \path let
  %   \p{a} = (rightNodeA.east),
  %   \p{b} = (rightNodeB.east),
  %   \n{rx} = {max(\x{a}, \x{b}) + \loopClearance}
  % in coordinate (R2) at (\n{rx}, 0);
  %
  % \path let
  %   \p{a} = (leftNodeA.west),
  %   \p{b} = (leftNodeB.west),
  %   \n{lx} = {min(\x{a}, \x{b}) - \loopClearance}
  % in coordinate (L2) at (\n{lx}, 0);
  %
  % Part 2: Enter target decision via degree anchors (.30 / .150)
  %         to avoid visual merging with side-exit arrows (.east/.west)
  %         and main-spine arrows (.north/.south).
  %
  % \draw [arrow]
  %   (rightNode.east) -- (R2 |- rightNode.east)   % horizontal to R2
  %                    -- (R2 |- target.30)         % vertical to .30 y
  %                    -- (target.30)               % horizontal into .30
  % ;
  %
  % \draw [arrow]
  %   (leftNode.west) -- (L2 |- leftNode.west)
  %                   -- (L2 |- target.150)
  %                   -- (target.150)
  % ;

\end{tikzpicture}

\end{document}
